\section{Conclusion}

\subsection{Summary of Findings}

This paper has made four primary claims about the Interstate 2.0 transit layer, each supported by the hub siting, marginal cost, city pair, and equity analyses presented above:

\textbf{C1}: Nine T1/T1 diamond hubs and approximately 50 T1/T2 regional stops bring an estimated 12.4 million transit-dependent Americans within 30 miles of an intercity bus connection they currently lack. The estimate is conservative: it uses the 30-mile radius rather than the 50-mile radius, applies a non-overlapping population assignment, and counts only zero-vehicle households and rural households without alternative transit, not all households that would benefit from improved intercity bus service.

\textbf{C2}: The hub increment investment --- \$2 billion on a \$253 billion program --- serves transit-dependent travelers at \$161 per traveler served: 14 times more cost-efficient than comparable standalone bus terminal construction (\$2,200 per traveler) and 17 times more efficient than Amtrak capital investment per rider (\$2,800 per rider). The efficiency advantage is structural, not accidental: it results from colocating passenger infrastructure with freight infrastructure that has already paid for the access roads, utilities, and parking.

\textbf{C3}: More than 40 origin-destination city pairs gain first-time intercity bus connectivity through the T1 hub network. The Southeast cluster (Atlanta, Jacksonville, Richmond, and T1/T2 regional stops) represents the largest single concentration of unmet demand; the transcontinental tier (Sacramento--Gary, Seattle--Gary) represents a set of routes that are currently infeasible due to PTI, but become schedulable at the managed freight lane's PTI of 1.15.

\textbf{C4}: T1/T1 hub locations align disproportionately with communities scoring high on C3 (Economic Opportunity Access), with Pearson $r = 0.68$ between C3 score and transit-dependent household density at hub locations. The rubric-driven freight site selection produces near-optimal transit equity siting as a byproduct, because the same economic geography that generates freight intensity also generates transit dependence.

\subsection{The Transit Layer as Design Requirement}

The central policy recommendation of this paper is that transit integration should be treated as a mandatory design requirement for all T1/T1 diamond hub construction, not as an optional add-on. The retrofit cost penalty for freight-only hub design is 3--5 times the incremental cost of transit-compatible design from the outset. Diamond hubs designed without transit-compatible geometry --- bus turning radii in access roads, ADA-accessible pathways from parking to platform, utility conduit routing for passenger facility connections --- cannot be efficiently upgraded for transit service once construction is complete.

This is not a soft recommendation: it is an investment protection argument for the \$4.5 billion diamond interchange upgrade program. If the 9 T1/T1 hubs are built without transit design standards, the \$2 billion transit layer cannot be effectively deployed at those locations, and the 14-fold cost efficiency advantage over standalone terminals is lost. The transit integration cost rises from 0.8\% to 3--4\% of total program cost if retrofit is required. The mandatory design standard is the less expensive option.

\subsection{Limitations}

The 12.4 million transit-dependent Americans figure is based on ACS B08201 zero-vehicle household counts within 30 miles of hub locations. Actual ridership will depend on factors not modeled here: bus operator market entry decisions (which depend on route economics, fleet availability, and regulatory environment), fare structure (which determines effective demand among the transit-dependent population), and service frequency (below which ridership does not develop because the service is not reliable enough for work or healthcare trips).

The city pair demand estimates are derived from a gravity model calibrated to pre-2020 Greyhound ridership data. The 2020--2026 period has seen significant changes in travel behavior --- remote work, reduced commuting, increased long-distance leisure travel --- that the pre-2020 calibration may not fully capture. The estimates should be treated as order-of-magnitude guidance for policy analysis, not as ridership projections for operator business plans.

\subsection{Next: F.2}

The companion F.2 paper computes travel times, transfer requirements, and equity metrics for specific T1 bus corridors in the Southeast cluster and the Gulf Coast alignment. F.2 moves from the hub siting and aggregate population analysis of this paper to the corridor-level service design: which routes are viable at what frequency, which populations they serve on each route, and what the equity outcomes are under different fare and frequency scenarios.

The transit layer is not an add-on to Interstate 2.0. It is the evidence that the highway investment produces larger social returns than the freight B/C ratio captures --- and that the country, having built the freight infrastructure, should not leave 12.4 million transit-dependent Americans on the wrong side of the platform.
